\section{Introduction}

Computational pathology is often divided into separate model problems: patch classification, foundation-model feature extraction, multiple-instance learning, domain robustness, or federated learning. PA-NF is built from the opposite direction. It treats the pathology system as one connected pipeline in which acquisition and institutional effects can enter at several stages and can be amplified as information moves from patches to representations, from representations to slide predictions, and from local models to collaborative models.

The name \emph{Paired-Acquisition Neural Factorization} originated in the representation stage, where exact same-tissue acquisitions expose what changes with scanner while holding the biological region fixed. That paired design remains the foundation of the program, but PA-NF now denotes the full pipeline built around it. Within this manuscript, \emph{PA-NF representation factorization} refers specifically to the tissue/acquisition two-branch model, while \emph{the PA-NF pipeline} refers to the complete computational pathology system.

The pipeline contains four principal authored systems:
\begin{itemize}
\item \textbf{PA-NF representation factorization:} a paired-scanner representation model with explicit tissue-oriented and acquisition-oriented branches, a joint decoder, crossed reconstruction, paired consistency, variance control, and scanner-routing objectives;
\item \textbf{TransnnMIL:} a whole-slide multiple-instance learning architecture family combining transformer-style global correlation modeling with gated aggregation and optional hierarchical, topological, and pruning components;
\item \textbf{PathologyFL:} a pathology-specific federated-learning framework with coordinator/client training, FedAvg/FedProx/FedAdam and custom weighting paths, robustness mechanisms, privacy integrations, monitoring, asynchronous execution, and failure handling;
\item \textbf{FAIR-WEIGHTS-H and dominance-aware aggregation:} institutional weighting and detector-switch mechanisms that test when sample-count authority should yield to cross-site evidence, uncertainty, and reliability signals.
\end{itemize}

These systems share a common experimental layer. PCam provides patch-level training and evaluation. SCORPION and multi-scanner canine SCC provide exact paired acquisitions across five scanners. PANDA provides 10,611 readable slide-level Phikon feature bags for whole-slide modeling and simulated institutional stress. CAMELYON17/WILDS supplies 455,954 examples across five centers for natural center-shift studies. Controlled latent-factor generators provide ground truth for mechanism questions that cannot be observed directly in pathology data.

The central scientific hypothesis is that scanner and site variation should not be treated as a single preprocessing nuisance. A scanner shortcut can be encoded in a patch representation, attended to by a slide model, and then amplified through institutional aggregation. Conversely, institutional reweighting can change which representation failures dominate the global model. PA-NF therefore studies the transitions between these levels rather than optimizing them in isolation.

The main contributions of this paper are: (i) a registered paired-acquisition representation result on SCORPION with an equal-capacity neural control and cross-backbone transfer; (ii) an independent multi-scanner canine study that tests the limits of the representation claim against strong simple corrections; (iii) an implemented whole-slide architecture and full-PANDA stabilization study; (iv) a custom federated pathology stack with a dominance-aware detector that improves over FedAvg in controlled PANDA failure regimes and transfers without retuning to systematic ordinal shift; and (v) a CAMELYON17 held-out-center study showing that source influence and source sample count are not interchangeable.

The goal is not to collapse these experiments into one leaderboard. Each dataset answers a different question at a different level of the pipeline. The result is a single research system with several experimentally linked components, not a claim that one metric or one model dominates all of computational pathology.
